\chapter{Wstęp}
\label{ch:wstep}

% Cel rozdziału: wprowadzić problem, cel, zakres i wkład pracy.
% Do napisania od zera (rozdział krótki, 3-5 stron).

\section{Motywacja}
\label{sec:motywacja}
% Problem weryfikacji dyplomów: fałszerstwa, koszt weryfikacji ręcznej,
% zależność od wydawcy (uczelnia musi żyć i odpowiadać na zapytania).
\TODO{Opisz problem weryfikacji dyplomów i jego znaczenie.}

\section{Cel i zakres pracy}
\label{sec:cel}
% Cel: zaprojektować i zmierzyć tani, prywatny rejestr dyplomów na Ethereum.
% Zakres: L1 + L2, prywatność (selektywne ujawnianie), zarządzanie.
\TODO{Sprecyzuj cel i granice zakresu.}

\section{Wkład własny}
\label{sec:wklad}
% Trzy główne wkłady (mapują się na rozdziały 5-7):
\begin{enumerate}
  \item Selektywne ujawnianie pól przez zobowiązania kryptograficzne
        (rozdz.~\ref{ch:prywatnosc}).
  \item Wdrożenie i pomiary kosztów na sieciach warstwy~2
        (rozdz.~\ref{ch:l2}).
  \item Model zarządzania (multisig + timelock) i analiza wartości
        blockchaina względem alternatyw (rozdz.~\ref{ch:governance}).
\end{enumerate}
\TODO{Rozwiń każdy wkład w 1-2 zdaniach.}

\section{Struktura pracy}
\label{sec:struktura}
\TODO{Krótki przewodnik po rozdziałach 2--8.}
